
\documentclass{article}
\usepackage{colortbl}
\usepackage{makecell}
\usepackage{multirow}
\usepackage{supertabular}

\begin{document}

\newcounter{utterance}

\centering \large Interaction Transcript for game `wordle', experiment `medium\_frequency\_words\_no\_clue\_no\_critic', episode 2 with nemotron{-}3{-}ultra.
\vspace{24pt}

{ \footnotesize  \setcounter{utterance}{1}
\setlength{\tabcolsep}{0pt}
\begin{supertabular}{c@{$\;$}|p{.15\linewidth}@{}p{.15\linewidth}p{.15\linewidth}p{.15\linewidth}p{.15\linewidth}p{.15\linewidth}}
    \# & \multicolumn{2}{c}{Player} && \multicolumn{2}{c}{Game Master} \\
    \hline

    \theutterance \stepcounter{utterance}  
    & & & \multicolumn{4}{p{0.6\linewidth}}{
        \cellcolor[rgb]{0.9,0.9,0.9}{
            \makecell[{{p{\linewidth}}}]{
                \texttt{\tiny{[P1$\langle$GM]}}
                \texttt{Int saħħar tal{-}lingwa li jħobb jaqta' kliem billi juża r{-}regoli mogħtija.} \\
\\ 
\texttt{Merħba f'Wordle! Għandek sitt tentattivi biex taqta' l{-}kelma fil{-}mira, kelma Maltija valida ta' ħames ittri żgħar. Jekk jogħġbok uża t{-}tikketti "spjegazzjoni:" u "tentattiv:" biex tipprovdi spjegazzjoni qasira għal kull tentattiv.} \\
\\ 
\texttt{Pereżempju, jekk it{-}tentattiv tiegħek huwa "darba", it{-}tweġiba tiegħek għandha tkun} \\
\texttt{spjegazzjoni: din hija kelma Maltija komuni ta' ħames ittri, u qed nibda t{-}tentattiv tiegħi b'din il{-}kelma.} \\
\texttt{tentattiv: darba} \\
\\ 
\texttt{Wara kull tentattiv, it{-}tweġiba tiegħek tiġi vvalidata, u tirċievi feedback li jindika liema ittri huma korretti (<aħdar>), liema ittri huma korretti iżda f'pożizzjoni ħażina (<isfar>), u liema ittri huma żbaljati (<ahmar>). Dan il{-}feedback jista' jkun utli biex tiddetermina liema ittri għandek tinkludi jew teskludi fit{-}tentattiv li jmiss.} \\
\\ 
\texttt{Pereżempju, il{-}feedback għal "darba" jista' jkun:} \\
\texttt{feedback tat{-}tentattiv: d<isfar> a<isfar> r<aħdar> b<isfar> a<ahmar>} \\
\\ 
\texttt{L{-}ispjegazzjoni għandha tinkludi dettalji dwar kif il{-}feedback tat{-}tentattiv jintuża biex tasal għal tentattiv ġdid.} \\
\\ 
\texttt{Tiġġenerax ebda test ieħor.} \\
\\ 
\texttt{Nibdew bl{-}ewwel tentattiv tiegħek.} \\
            }
        }
    }
     \\ \\

    \theutterance \stepcounter{utterance}  
    & \multicolumn{4}{p{0.6\linewidth}}{
        \cellcolor[rgb]{0.561,0.737,0.561}{
            \makecell[{{p{\linewidth}}}]{
                \texttt{\tiny{[P1$\rangle$GM]}}
                \texttt{spjegazzjoni: Nibda b'kelma komuni li tieħu ħafna vokaali (a, e, i) u konsonanti komuni (r, t, l) biex nemmennu jew nestkludu l{-}ittri l{-}akbar numru ta' volte.} \\
\texttt{tentattiv: rieti} \\
            }
        }
    }
    & & \\ \\

    \theutterance \stepcounter{utterance}  
    & & & \multicolumn{2}{p{0.3\linewidth}}{
        \cellcolor[rgb]{0.9,0.9,0.9}{
            \makecell[{{p{\linewidth}}}]{
                \texttt{\tiny{[GM$|$GM]}}
                \texttt{Error: The guessed word is not a valid word for this game.} \\
            }
        }
    }
    & & \\ \\

    \theutterance \stepcounter{utterance}  
    & & & \multicolumn{4}{p{0.6\linewidth}}{
        \cellcolor[rgb]{0.9,0.9,0.9}{
            \makecell[{{p{\linewidth}}}]{
                \texttt{\tiny{[P1$\langle$GM]}}
                \texttt{It{-}tentattiv tiegħek mhuwiex kelma valida għal din il{-}logħba. Jekk jogħġbok erġa' pprova.} \\
\\ 
\texttt{Ipprovdi t{-}tweġiba tiegħek biss f'dan il{-}format.} \\
\texttt{spjegazzjoni: dettalji} \\
\texttt{tentattiv: kelma} \\
            }
        }
    }
     \\ \\

    \theutterance \stepcounter{utterance}  
    & \multicolumn{4}{p{0.6\linewidth}}{
        \cellcolor[rgb]{0.561,0.737,0.561}{
            \makecell[{{p{\linewidth}}}]{
                \texttt{\tiny{[P1$\rangle$GM]}}
                \texttt{spjegazzjoni: "rieti" mhux kelma valida f'dan il{-}kontext (forse mhux f'word{-}list ta' Wordle Malti). Nibdel strategija u nuzzaw "qatel" li hija kelma verba Maltija komuni (forma imperattiv/qies), tieħu l{-}istess vokaali (a, e) u r{-}r u t li huma konsonanti ħafna komuni, u t{-}tlieta it{-}tri differenti mill{-}ewwel.} \\
\texttt{tentattiv: qatel} \\
            }
        }
    }
    & & \\ \\

    \theutterance \stepcounter{utterance}  
    & & & \multicolumn{2}{p{0.3\linewidth}}{
        \cellcolor[rgb]{0.9,0.9,0.9}{
            \makecell[{{p{\linewidth}}}]{
                \texttt{\tiny{[GM$|$GM]}}
                \texttt{attempts = 1} \\
\texttt{target\_word = ħajtu} \\
\texttt{guess = qatel} \\
\texttt{guess\_feedback = q<red> a<green> t<yellow> e<red> l<red>} \\
            }
        }
    }
    & & \\ \\

    \theutterance \stepcounter{utterance}  
    & & & \multicolumn{4}{p{0.6\linewidth}}{
        \cellcolor[rgb]{0.9,0.9,0.9}{
            \makecell[{{p{\linewidth}}}]{
                \texttt{\tiny{[P1$\langle$GM]}}
                \texttt{feedback tat{-}tentattiv: q<ahmar> a<aħdar> t<isfar> e<ahmar> l<ahmar>} \\
\\ 
\texttt{Ipprovdi t{-}tweġiba tiegħek biss f'dan il{-}format.} \\
\texttt{spjegazzjoni: dettalji} \\
\texttt{tentattiv: kelma} \\
            }
        }
    }
     \\ \\

    \theutterance \stepcounter{utterance}  
    & \multicolumn{4}{p{0.6\linewidth}}{
        \cellcolor[rgb]{0.561,0.737,0.561}{
            \makecell[{{p{\linewidth}}}]{
                \texttt{\tiny{[P1$\rangle$GM]}}
                \texttt{spjegazzjoni: L{-}ittra 'a' hija f'pożizzjoni korretta (it{-}tieni). L{-}ittra 't' jokkupa fid{-}darba (ikun it{-}tielet jew ir{-}raba' ittra). L{-}ittr 'q', 'e', u 'l' mhumiex preżenti fil{-}kelma. Niktbu 't' fit{-}tielet pożizzjoni u niftakar li l{-}ewwel u r{-}raba' ittra għandhom jkunu minn: i, o, u, b, d, f, g, h, j, k, m, n, p, r, s, v, w, x, z, ċ, ġ, għ, ħ, ż. Nnazzlu 'matir' (kelma Maltija li tiena 'matrix' jew 'mould').} \\
\texttt{tentattiv: matir} \\
            }
        }
    }
    & & \\ \\

    \theutterance \stepcounter{utterance}  
    & & & \multicolumn{2}{p{0.3\linewidth}}{
        \cellcolor[rgb]{0.9,0.9,0.9}{
            \makecell[{{p{\linewidth}}}]{
                \texttt{\tiny{[GM$|$GM]}}
                \texttt{Error: The guessed word is not a valid word for this game.} \\
            }
        }
    }
    & & \\ \\

    \theutterance \stepcounter{utterance}  
    & & & \multicolumn{4}{p{0.6\linewidth}}{
        \cellcolor[rgb]{0.9,0.9,0.9}{
            \makecell[{{p{\linewidth}}}]{
                \texttt{\tiny{[P1$\langle$GM]}}
                \texttt{It{-}tentattiv tiegħek mhuwiex kelma valida għal din il{-}logħba. Jekk jogħġbok erġa' pprova.} \\
\\ 
\texttt{Ipprovdi t{-}tweġiba tiegħek biss f'dan il{-}format.} \\
\texttt{spjegazzjoni: dettalji} \\
\texttt{tentattiv: kelma} \\
            }
        }
    }
     \\ \\

    \theutterance \stepcounter{utterance}  
    & \multicolumn{4}{p{0.6\linewidth}}{
        \cellcolor[rgb]{0.561,0.737,0.561}{
            \makecell[{{p{\linewidth}}}]{
                \texttt{\tiny{[P1$\rangle$GM]}}
                \texttt{spjegazzjoni: "matir" mhux kelma valida f'dan il{-}kontext. Nieħdu 't' fix{-}xewk pożizzjoni (daqs li l{-}feedback kien isfar għall{-}it{-}tieni pożizzjoni, u l{-}ewwel u t{-}tiela huma mblokka minn q, e, l). Niktbu 'b' u 'n' bħala konsonanti komuni fl{-}ewwel u r{-}raba' pożizzjoni, u 'u' bħala vokaali ġdida fit{-}tlieta pożizzjoni. Kelma valida: 'bantuh' mhux 5 ittri. 'bantu' (plural ta' ban/band). Niprova 'bantu'.} \\
\texttt{tentattiv: bantu} \\
            }
        }
    }
    & & \\ \\

    \theutterance \stepcounter{utterance}  
    & & & \multicolumn{2}{p{0.3\linewidth}}{
        \cellcolor[rgb]{0.9,0.9,0.9}{
            \makecell[{{p{\linewidth}}}]{
                \texttt{\tiny{[GM$|$GM]}}
                \texttt{Error: The guessed word is not a valid word for this game.} \\
            }
        }
    }
    & & \\ \\

    \theutterance \stepcounter{utterance}  
    & & & \multicolumn{4}{p{0.6\linewidth}}{
        \cellcolor[rgb]{0.9,0.9,0.9}{
            \makecell[{{p{\linewidth}}}]{
                \texttt{\tiny{[P1$\langle$GM]}}
                \texttt{It{-}tentattiv tiegħek mhuwiex kelma valida għal din il{-}logħba. Jekk jogħġbok erġa' pprova.} \\
\\ 
\texttt{Ipprovdi t{-}tweġiba tiegħek biss f'dan il{-}format.} \\
\texttt{spjegazzjoni: dettalji} \\
\texttt{tentattiv: kelma} \\
            }
        }
    }
     \\ \\

    \theutterance \stepcounter{utterance}  
    & \multicolumn{4}{p{0.6\linewidth}}{
        \cellcolor[rgb]{0.561,0.737,0.561}{
            \makecell[{{p{\linewidth}}}]{
                \texttt{\tiny{[P1$\rangle$GM]}}
                \texttt{spjegazzjoni: "bantu" mhix kelma valida. Nibdew minn 't' f'pożizzjoni 3. Nprova kombinazzjoni b'vokaali 'i' u 'u' li huma komuni. Niktbu 'fit' (from 'feta')? Le, 'f' ahmar. 'kit'? Le. 'sit'? Le. 'hit'? Le. 'bit'? Le. 'nit'? Le. 'git'? Le. 'lit'? Le. 'mit'? Le. 'pit'? Le. 'rit'? Le. 'sit'? Le. 'tit'? Le. 'vit'? Le. 'xit'? Le. 'zit'? Le. 'ċit'? Le. 'ġit'? Le. 'għit'? Le. 'ħit'? Le. 'żit'? Le. 'tift'? Le. 'tibt'? Le. 'tikt'? Le. 'tilt'? Le. 'tint'? Le. 'tist'? Le. 'titf'? Le. 'titk'? Le. 'titm'? Le. 'titn'? Le. 'titp'? Le. 'titr'? Le. 'tits'? Le. 'titt'? Le. 'titx'? Le. 'titz'? Le. 'titċ'? Le. 'titġ'? Le. 'titgħ'? Le. 'titħ'? Le. 'titż'? Le. } \\
\texttt{Issa nprova 'a' f'pożizzjoni 1, 't' f'3, 'u' f'5 {-}> 'atu' qisar. 'matu'? 'patu'? 'ratu'? 'satu'? 'batu'? 'datu'? 'fatu'? 'gatu'? 'hatu'? 'jatu'? 'katu'? 'latu'? 'natu'? 'vatu'? 'watu'? 'xatu'? 'zatu'? 'ċatu'? '} \\
            }
        }
    }
    & & \\ \\

    \theutterance \stepcounter{utterance}  
    & & & \multicolumn{2}{p{0.3\linewidth}}{
        \cellcolor[rgb]{0.9,0.9,0.9}{
            \makecell[{{p{\linewidth}}}]{
                \texttt{\tiny{[GM$|$GM]}}
                \texttt{Error: The response should contain the 'tentattiv:' keyword exactly once.} \\
            }
        }
    }
    & & \\ \\

    \theutterance \stepcounter{utterance}  
    & & & \multicolumn{2}{p{0.3\linewidth}}{
        \cellcolor[rgb]{0.9,0.9,0.9}{
            \makecell[{{p{\linewidth}}}]{
                \texttt{\tiny{[GM$|$GM]}}
                \texttt{game\_result = ABORT} \\
            }
        }
    }
    & & \\ \\

\end{supertabular}
}

\end{document}
